%% kapa: exam=1 round=20 year=2009 period=1 form=A subject=civil_law subject_ko=민법 seq=7 number=7
%% source: Q-Net 감정평가사 기출문제 (cst00309) artlSeq=5024178 file=extracted/1008/5024178/2009년 제20회 감정평가사 1차시험 기출문제/1교시.hwp
%% conversion: kordoc 4.24.51 -> tools/gen1.py (auto); review=proofread-sonnet; flags=none
\begin{question}{7}
법원은 부재자 甲의 재산관리인으로 乙을 선임하였다. 그 후 乙은 법원의 허가를 얻어 甲의 토지를 丙에게 매도ㆍ등기하였다. 다음 중 옳지 \uline{않은} 것은? (다툼이 있으면 판례에 의함)
\begin{choices}
  \item 乙은 원칙적으로 甲의 재산의 관리행위를 할 권한을 가진다.
  \item 乙이 보존행위와 처분행위를 하는 경우에 법원의 허가가 필요하다.
  \item 甲이 생환하여 재산관리인 乙의 선임결정의 취소를 청구한 경우 법원은 이를 취소하여야 한다.
  \item ③의 경우 선임결정의 취소는 소급효가 없으므로 乙과 丙의 토지매매계약은 유효하다.
  \item 만일 甲의 사망이 확인되었더라도 乙의 권한이 당연히 소멸하는 것은 아니다.
\end{choices}
\end{question}
\answer{2}
